\documentclass[11pt,a4paper]{article}
\usepackage[margin=1in]{geometry}
\usepackage[T1]{fontenc}
\usepackage[utf8]{inputenc}
\usepackage{lmodern}
\usepackage{microtype}
\usepackage{amsmath,amssymb,mathtools}
\usepackage[hidelinks]{hyperref}
\usepackage{fancyhdr}
\setlength{\headheight}{14pt}
\pagestyle{fancy}
\fancyhf{}
\lhead{Ernst Steinitz}
\rhead{Deutscher Text}
\cfoot{\thepage}
\setlength{\parindent}{1.5em}
\setlength{\parskip}{0.35em}
\allowdisplaybreaks
\newcommand{\Cl}{\operatorname{Cl}}
\newcommand{\Md}{\operatorname{Md}}
\newcommand{\Nm}{\operatorname{Nm}}
\newcommand{\Rg}{\operatorname{Rg}}
\newcommand{\Gd}{\operatorname{Gd}}
\newcommand{\Dg}{\operatorname{Dg}}
\newcommand{\tht}{\vartheta}
\begin{document}
\section*{Zur Theorie der Moduln}
\begin{center}Von \emph{Ernst Steinitz} in Charlottenburg.\end{center}

Es seien $A_1,\ldots,A_s,\Gamma,\Delta$ beliebige Classen $n$ten Grades, und
\[
 \mathfrak a^{(i)}_1,\mathfrak a^{(i)}_2,\ldots,
 \mathfrak a^{(i)}_{\psi(A_i)}\qquad (i=1,2,\ldots,s)
\]
die sämmtlichen Moduln von $A_i$. Die Anzahl der Systeme zu je $s$ Moduln
\begin{equation}
 \mathfrak a^{(1)}_{q_1},\mathfrak a^{(2)}_{q_2},\ldots,
 \mathfrak a^{(s)}_{q_s},
\tag{1}
\end{equation}
welche bez. den Classen $A_1,A_2,\ldots,A_s$ angehören, beträgt
\[
 \psi(A_1)\psi(A_2)\cdots\psi(A_s).
\]
Die Anzahl der Systeme (1), für welche das kleinste gemeinsame Vielfache
\begin{equation}
 [\mathfrak a^{(1)}_{q_1},\mathfrak a^{(2)}_{q_2},\ldots,
 \mathfrak a^{(s)}_{q_s}]
\tag{2}
\end{equation}
gleich einem bestimmten Modul einer Classe $\Theta$ wird, ist durch die Function
$\bar\gamma(\Theta;A_1,\ldots,A_s)$ dargestellt. Ist daher $\mathfrak c$ irgend ein Modul von $\Gamma$, so ist die Anzahl derjenigen Systeme (1), für welche (2) irgend ein der Classe $\Theta$ angehöriger Divisor von $\mathfrak c$ wird, gleich
\[
 t(\Gamma\mid\Theta)\bar\gamma(\Theta;A_1,\ldots,A_s).
\]
Um endlich die Anzahl aller Systeme (1) zu erhalten, für welche (2) irgend ein Divisor von $\mathfrak c$ wird, hat man nur die Summe
\[
 \sum_{\Theta}t(\Gamma\mid\Theta)\bar\gamma(\Theta;A_1,\ldots,A_s)
\]
zu bilden, in welcher $\Theta$ alle Classen $n$ten Grades durchläuft. Bedenkt man aber, dass der Modul (2) stets und nur dann Divisor von $\mathfrak c$ wird, wenn alle Moduln des Systems (1) Divisoren von $\mathfrak c$ sind, so erkennt man, dass die zuletzt berechnete Anzahl auch durch den Ausdruck
\[
 t(\Gamma\mid A_1)t(\Gamma\mid A_2)\cdots t(\Gamma\mid A_s)
\]
dargestellt wird. Man erhält daher die Relation
\[
\text{I.}\qquad
 \sum_{\Theta}t(\Gamma\mid\Theta)\bar\gamma(\Theta;A_1,\ldots,A_s)
 =t(\Gamma\mid A_1)t(\Gamma\mid A_2)\cdots t(\Gamma\mid A_s).
\]
Diese für jedes System von Classen $\Gamma,A_1,\ldots,A_s$ geltende Relation reicht nun vollkommen hin, um die Function $\bar\gamma$ zu definiren.

Zunächst erkennt man aus I, dass $\bar\gamma(\Gamma;A_1,\ldots,A_s)$ verschwindet, wenn $\Gamma$ nicht Vielfaches jeder der Classen $A_1,\ldots,A_s$ ist. In diesem Falle nämlich verschwindet das Product auf der rechten Seite von I. Setzen wir jetzt unsere Behauptung für alle Divisoren von $\Gamma$ mit Ausnahme von $\Gamma$ selbst als bewiesen voraus, so verschwindet für jedes von $\Gamma$ verschiedene $\Theta$ ein Factor des Productes unter dem Summenzeichen, während $t(\Gamma\mid\Gamma)=1$ wird. Man erhält daher
\[
 \bar\gamma(\Gamma;A_1,\ldots,A_s)=0.
\]
Hieraus folgt, dass man für jede Classe $\Gamma$ der Relation I die Form geben kann:
\begin{equation}
 \sum_{\substack{(\Theta,\Gamma)=\Theta\\(\Theta,\Gamma')=\Gamma'}}
 t(\Gamma\mid\Theta)\bar\gamma(\Theta;A_1,\ldots,A_s)
 =t(\Gamma\mid A_1)\cdots t(\Gamma\mid A_s),
\tag{3}
\end{equation}
wo $\Gamma'$ das kleinste gemeinsame Vielfache der Classen $A_1,\ldots,A_s$ bezeichnet und die Summation sich auf alle zwischen $\Gamma$ und $\Gamma'$ (mit Einschluss der Grenzen) liegende Classen $\Theta$ erstreckt. Ist daher $\Gamma=\Gamma'$, so erhält man einfach
\[
 \bar\gamma(\Gamma';A_1,\ldots,A_s)=t(\Gamma'\mid A_1)\cdots t(\Gamma'\mid A_s),
\]
ist aber $\Gamma$ ein Vielfaches von $\Gamma'$, jedoch von $\Gamma'$ verschieden, so denken wir uns der Berechnung von $\bar\gamma(\Gamma;A_1,\ldots,A_s)$ die Berechnung derjenigen nur in endlicher Anzahl vorhandenen Ausdrücke $\bar\gamma(\Theta;A_1,\ldots,A_s)$ vorangegangen, in denen $\Theta$ zwischen $\Gamma$ und $\Gamma'$ liegt und von $\Gamma$ verschieden ist. Man erhält dann aus (3)
\[
\text{Ia.}\quad
 \bar\gamma(\Gamma;A_1,\ldots,A_s)
 =t(\Gamma\mid A_1)\cdots t(\Gamma\mid A_s)
 -\sum_{\substack{(\Theta,\Gamma)=\Theta\\(\Theta,\Gamma')=\Gamma'\\\Theta\ne\Gamma}}
 t(\Gamma\mid\Theta)\bar\gamma(\Theta;A_1,\ldots,A_s),
\]
und hier stehen auf der rechten Seite nur bekannte Grössen.

In gleicher Weise erhält man für $\bar\delta$, wenn
\[
 \Delta'=(A_1,\ldots,A_s)
\]
gesetzt wird:
\[
\text{II.}\qquad
 \sum_{\Theta}v(\Theta\mid\Delta)\bar\delta(\Theta;A_1,\ldots,A_s)
 =v(A_1\mid\Delta)\cdots v(A_s\mid\Delta),
\]
\[
\text{IIa.}\quad
 \bar\delta(\Delta;A_1,\ldots,A_s)=v(A_1\mid\Delta)\cdots v(A_s\mid\Delta)
 -\sum_{\substack{(\Theta,\Delta)=\Delta\\(\Theta,\Delta')=\Theta\\\Theta\ne\Delta}}
 v(\Theta\mid\Delta)\bar\delta(\Theta;A_1,\ldots,A_s).
\]
Die Recursionsformeln Ia, IIa zeigen, dass man die Functionen $\bar\gamma$ und $\bar\delta$ aus $t$- bez. $v$-Functionen durch die Operationen Addition, Subtraction, Multiplication zusammensetzen kann. Diese Darstellung ist aber wesentlich complicirter als die Darstellung der Functionen $t$ und $v$ durch $\psi$-Functionen. Während nämlich bei dieser die Anzahl der darstellenden $\psi$-Functionen durch den Grad $n$ der betrachteten Classen beschränkt ist, ist die Anzahl der bei der Darstellung der Functionen $\bar\gamma(\Gamma;A_1,\ldots,A_s)$, $\bar\delta(\Delta;A_1,\ldots,A_s)$ verwandten $t$- und $v$-Functionen zwar in jedem einzelnen Falle eine endliche, sie ist aber, wenn auch der Grad $n$ und die Zahl $s$ gegeben sind, an keine obere Grenze gebunden. Ferner wird durch die hier in den Formeln auftretende Subtraction die Beantwortung der Frage nach dem Verschwinden der $\bar\gamma$- und $\bar\delta$-Function sehr erschwert. Diese Schwierigkeiten werden beseitigt, indem es gelingt, die für die Functionen $\bar\gamma$ und $\bar\delta$ gewonnenen Ausdrücke umzuformen in Producte von Functionen, welche eine leichtere Übersicht gestatten.

\section*{§ 10. Vorbereitende Betrachtungen.}

1. Wir beschränken uns von nun an auf die Betrachtung primärer Classen. Eine solche Classe
\[
 A=\Cl(p^{r_1}\mid p^{r_2}\mid\cdots\mid p^{r_n})
\]
ist vollkommen bestimmt durch das Exponentensystem
\begin{equation}
 r_1,r_2,\ldots,r_n
\tag{1}
\end{equation}
und die Primzahl $p$. Denken wir uns nur das Exponentensystem (1) gegeben, so haben wir es mit einer sog. ,,halbbestimmten`` primären Classe zu thun. Unter den kleinen griechischen Buchstaben sind im Folgenden stets halbbestimmte primäre Classen zu verstehen. Ist $\alpha$ eine solche Classe und (1) das System ihrer Exponenten, so schreiben wir auch
\[
 \alpha=\Cl(r_1\mid r_2\mid\cdots\mid r_n).
\]
Ist wie in § 6
\[
\begin{gathered}
 r_1=\cdots=r_{s_1}=e_1,
 \quad r_{s_1+1}=\cdots=r_{s_1+s_2}=e_2,
 \quad \ldots,\quad
 r_{n-s_k+1}=\cdots=r_n=e_k,\
 e_1<e_2<\cdots<e_k,
\end{gathered}
\]
so erhält man (§ 6, III; a, b)
\begin{align}
\psi(\alpha)
&=p^{\sum_{i>h}s_i s_h(e_i-e_h)}
 {\vartheta(p;n)\over \vartheta(p;s_1)\vartheta(p;s_2)\cdots\vartheta(p;s_k)} \notag\\
&=p^{\sum_{i>h}s_i s_h(e_i-e_h-1)}
 g(p;n,s_1)g(p;n-s_1,s_2)\cdots
 g(p;n-s_1-\cdots-s_{k-2},s_{k-1})
\tag{2}
\end{align}
als ganze Function der unbestimmten Primzahl $p$.

Werden mehrere Classen $\alpha,\beta,\ldots$ gleichzeitig betrachtet, so ist vorausgesetzt, dass für alle die unbestimmte Primzahl $p$ dieselbe ist. Hiernach ist klar, was Ausdrücke wie ,,$\alpha$ ist theilbar durch $\beta$`` u. a. bedeuten. Hat man $s$ Classen
\[
 \alpha_i=\Cl(a_{i1}\mid a_{i2}\mid\cdots\mid a_{in})\qquad (i=1,\ldots,s),
\]
und ist
\[
 \gamma'=[\alpha_1,\ldots,\alpha_s]=\Cl(c'_1\mid\cdots\mid c'_n)
\]
ihr kleinstes gemeinsames Vielfache,
\[
 \delta'=(\alpha_1,\ldots,\alpha_s)=\Cl(d'_1\mid\cdots\mid d'_n)
\]
ihr grösster gemeinsamer Theiler, so ist $c'_k$ die grösste, $d'_k$ die kleinste der Zahlen $a_{1k},\ldots,a_{sk}$. Ferner ergiebt sich aus der in § 7 gegebenen Definition leicht, wie man das Grenzvielfache
\[
 \gamma''=[|\alpha_1,\ldots,\alpha_s|]
\]
und den Grenztheiler
\[
 \delta''=(|\alpha_1,\ldots,\alpha_s|)
\]
der Classen $\alpha_1,\ldots,\alpha_s$ zu ermitteln hat: Man ordne die $ns$ Exponenten von $\alpha_1,\ldots,\alpha_s$ nach ihrer Grösse, sodass sie eine nirgends abnehmende Reihe bilden. Sind $d''_1,\ldots,d''_n$ die $n$ kleinsten, $c''_1,\ldots,c''_n$ die $n$ grössten, so ist
\[
 \gamma''=\Cl(c''_1\mid\cdots\mid c''_n),\qquad
 \delta''=\Cl(d''_1\mid\cdots\mid d''_n).
\]
Dass $\alpha$ durch $\beta$ theilbar ist, drücken wir durch die Ungleichung $\alpha\ge\beta$ oder $\beta\le\alpha$ aus; die Bezeichnungen $\alpha>\beta$, $\beta<\alpha$ schliessen die Gleichheit der Classen $\alpha$ und $\beta$ aus.

Ist die Ungleichung $\alpha\ge\beta$ nicht erfüllt, so verschwinden $t(\alpha\mid\beta)$ und $v(\alpha\mid\beta)$ identisch, d. h. für jedes $p$. Sonst ergiebt sich die Bedeutung von $t(\alpha\mid\beta)$ und $v(\alpha\mid\beta)$ aus den Formeln § 8, III, IV, ebenso wie die der Functionen $\bar\gamma(\gamma;\alpha_1,\ldots,\alpha_s)$, $\bar\delta(\delta;\alpha_1,\ldots,\alpha_s)$ aus den Formeln § 9, Ia, IIa. Alle diese Functionen sind, wie man leicht erkennt, ganze Functionen von $p$. Die Darstellung gewinnt an Anschaulichkeit, wenn wir hierbei $p$ als eine im reellen Gebiete beliebig veränderliche Grösse betrachten. Wir bringen dies äusserlich zum Ausdruck, indem wir den Buchstaben $p$ durch $x$ ersetzen.

2. Somit haben wir jetzt
\[
 \psi(\alpha)=x^{\sum_{i>h}s_i s_h(e_i-e_h)}
 {\vartheta(x;n)\over \vartheta(x;s_1)\cdots\vartheta(x;s_k)}
 =x^{\sum_{i>h}s_i s_h(e_i-e_h-1)}
 g(x;n,s_1)\cdots g(x;n-s_1-\cdots-s_{k-2},s_{k-1}),
\]
wo $s_1,\ldots,s_k,e_1,\ldots,e_k$ ihre frühere Bedeutung haben. Sind also die Exponenten von $\alpha$ nicht numerisch gegeben, so muss man, um den Ausdruck $\psi(\alpha)$ bilden zu können, noch die Zahlen $s_1,\ldots,s_k$ haben, d. h. es muss von zwei aufeinanderfolgenden Exponenten noch gesagt sein, ob sie gleich sind oder nicht. Noch grösseren Unbequemlichkeiten ist man bei den Functionen $t$ und $v$ ausgesetzt, und die Behandlung der Functionen $\bar\gamma$ und $\bar\delta$ wird dadurch ganz unübersichtlich. Man kann diese Schwierigkeiten beseitigen, indem man die Exponenten durch andere Grössen ersetzt, welche unter dem Namen ,,Indices`` in die Betrachtung eingeführt werden mögen.

Als ,,Reihe der Indices`` einer Classe $\alpha$ bezeichnen wir eine unendliche Reihe ganzer Zahlen
\[
 a_1,a_2,a_3,\ldots,
\]
deren $h$tes Element $a_h$ angiebt, wie viele Exponenten von $\alpha$ kleiner als $h$ sind. So gehören z. B. zu den Classen
\[
 \Cl(2\mid3\mid3\mid5\mid8),\qquad \Cl(0\mid4\mid7)
\]
die Reihen
\[
 0,0,1,3,3,4,4,4,5,5,5,\ldots;
 \qquad
 1,1,1,1,2,2,2,3,3,\ldots.
\]
Man erkennt leicht die Richtigkeit der folgenden Sätze:

1) Die Reihe der Indices einer Classe $n$ten Grades ist eine unendliche niemals abnehmende Reihe ganzer nicht negativer Zahlen, welche bis zur Zahl $n$ ansteigt.

2) Umgekehrt ist jede Reihe, welche die unter 1) angegebenen Eigenschaften besitzt, Reihe der Indices für eine gewisse Classe $\alpha$ vom Grade $n$.

3) Der Definition zufolge giebt der $h$te Index von $\alpha$ an, wie viele Exponenten $<h$ sind.

4) Umgekehrt giebt der $h$te Exponent von $\alpha$ an, wie viele Indices $\le h$ sind.

Man kann sich die Reihe der Indices $a_1,a_2,\ldots$ einer Classe auch nach der negativen Seite hin fortgesetzt denken. Es ist dann $a_0=0$, $a_{-1}=0$, $a_{-2}=0,\ldots$ zu setzen.

Wir fügen nun zu den Sätzen 1) bis 4) noch die folgenden, deren Beweis ebenfalls leicht ist.

5) Sind $\alpha$ und $\beta$ Classen vom Grade $n$, $a_h$ die Indices von $\alpha$, $b_h$ die Indices von $\beta$, so ist die Ungleichung $\alpha\ge\beta$ dann und nur dann erfüllt, wenn für jedes $h$ $a_h\le b_h$ ist.

Es seien jetzt $\alpha_1,\alpha_2,\ldots,\alpha_s$ Classen $n$ten Grades und es werde
\[
 [\alpha_1,\alpha_2,\ldots,\alpha_s]=\gamma',\quad
 (\alpha_1,\alpha_2,\ldots,\alpha_s)=\delta',
\]
\[
 [|\alpha_1,\alpha_2,\ldots,\alpha_s|]=\gamma'',\quad
 (|\alpha_1,\alpha_2,\ldots,\alpha_s|)=\delta''
\]
gesetzt; mit $a_h^{(i)}$ $(i=1,2,\ldots,s)$, $c'_h,d'_h,c''_h,d''_h$ seien die Indices der Classen $\alpha_i$, $\gamma'$, $\delta'$, $\gamma''$, $\delta''$ bezeichnet. Aus 5) folgt:

6) Der $h$te Index $c'_h$ des kleinsten gemeinsamen Vielfachen $\gamma'$ von $\alpha_1,\ldots,\alpha_s$ ist gleich dem kleinsten, der $h$te Index $d'_h$ des grössten gemeinsamen Theilers $\delta'$ von $\alpha_1,\ldots,\alpha_s$ gleich dem grössten unter den $h$ten Indices $a_h^{(1)},\ldots,a_h^{(s)}$ von $\alpha_1,\ldots,\alpha_s$.

Für die Reihen der Indices des Grenzvielfachen und des Grenztheilers erhält man den Satz:

7) Der $h$te Index $c''_h$ des Grenzvielfachen $\gamma''$ der Classen $\alpha_1,\ldots,\alpha_s$ ist gleich der grösseren der beiden Zahlen
\[
 0\quad\hbox{und}\quad n-\sum_{i=1}^s(n-a_h^{(i)})=\biggl(\sum_{i=1}^s a_h^{(i)}\biggr)-(s-1)n,
\]
der $h$te Index $d''_h$ ihres Grenztheilers $\delta''$ gleich der kleineren der beiden Zahlen
\[
 n\quad\hbox{und}\quad \sum_{i=1}^s a_h^{(i)}.
\]
Dabei sind mit $a_h^{(i)}$ die Indices von $\alpha_i$ bezeichnet worden.

Aus 4) ergiebt sich leicht:

8) Zwischen den Indices $a_h$ und den Exponenten $\bar a_1,\ldots,\bar a_n$ einer Classe $\alpha$ besteht unter anderen die Relation
\begin{equation}
 (1-n)\bar a_1+(3-n)\bar a_2+(5-n)\bar a_3+\cdots+(n-1)\bar a_n
 =\sum_h a_h(n-a_h).
\tag{3}
\end{equation}
Die Grenzen, bis zu welchen die Summation erstreckt werden muss, brauchen nicht angegeben zu werden, da die Relation jedenfalls richtig ist, wenn man $h$ alle ganzen Zahlen durchlaufen lässt.

3. In § 6 ist die Function $g(x;n,m)$ eingeführt und für alle ganzzahligen Werthe von $n$ und $m$ definirt worden. Wir stellen hier einige Formeln und Sätze, welche diese Function betreffen und welche wir später brauchen, zusammen.
\[
\begin{array}{ll}
1)&\displaystyle g(x;n,m)={ (1-x^n)(1-x^{n-1})\cdots(1-x^{n-m+1})\over(1-x)(1-x^2)\cdots(1-x^m)}\quad(m>0),\\[1.0ex]
2)&g(x;n,0)=1,\\
3)&g(x;n,m)=0\quad(m<0),\\
4)&g(x;0,m)=0\quad(m\ne0),\qquad g(x;0,0)=1,\\
5)&g(x;n,m)=g(x;n-1,m)+x^{n-m}g(x;n-1,m-1),\\
6)&g(x;n,m)=g(x;n-1,m-1)+x^m g(x;n-1,m),\\
7)&g(x;n,m)=x^{m(n-m)}g(1/x;n,m),\\
8)&g(x;s,q)g(x;s-q,r)=g(x;s,r)g(x;s-r,q),\\
9)&g(x;-n,m)=(-1)^m x^{-mn-{m(m-1)\over2}}g(x;n+m-1,m),\\
10)&g(x;n,m)=g(x;n,n-m)\quad(n\ge0),\\[0.5ex]
11)&\displaystyle g(1/x;n,m)={\vartheta(x;n)\over\vartheta(x;m)\vartheta(x;n-m)}\quad(n\ge m\ge0).
\end{array}
\]
12) In allen Fällen giebt die Entwickelung von $g$ nach Potenzen von $x$ nur eine endliche Anzahl positiver und negativer Potenzen, und die Coefficienten sind ganze nicht negative Zahlen.

13) Für $n\ge0$ ist $g(x;n,m)$ eine ganze Function von $x$; für $n>0$, $m\ge0$ verschwindet dieselbe nicht identisch.

\section*{§ 11. Die Functionen $\psi(\alpha)$, $t(\alpha\mid\beta)$, $v(\alpha\mid\beta)$, dargestellt mit Hülfe der Indices.}

Hat die Classe
\[
 \alpha=\Cl(\bar a_1\mid\bar a_2\mid\cdots\mid\bar a_n)
\]
die Indices $a_h$, so stellt $a_h-a_{h-1}$ die Anzahl derjenigen Exponenten dar, welche gleich $h-1$ sind. Berücksichtigt man daher, dass $\vartheta(x;0)=1$ ist, so erhält man nach § 6, III die Formel
\begin{equation}
 \psi(\alpha)=x^{(1-n)\bar a_1+(3-n)\bar a_2+\cdots+(n-1)\bar a_n}
 {\vartheta(x;n)\over\prod_h\vartheta(x;a_h-a_{h-1})}.
\tag{1}
\end{equation}
Nun ist aber (§ 10, 2, 8))
\[
 (1-n)\bar a_1+(3-n)\bar a_2+\cdots+(n-1)\bar a_n=\sum_h(n-a_h)a_h,
\]
ferner
\begin{align*}
 {\vartheta(x;n)\over\prod_h\vartheta(x;a_h-a_{h-1})}
 &=\prod_h {\vartheta(x;n-a_{h-1})\over\vartheta(x;n-a_h)\vartheta(x;a_h-a_{h-1})}\\
 &=\prod_h g(1/x;n-a_{h-1},n-a_h)\\
 &=\prod_h g(1/x;a_h,a_{h-1}).
\end{align*}
Daher erhalten wir unter Berücksichtigung von § 10, 3, 7)
\begin{align*}
\text{I.}\quad
\psi(\alpha)&=x^{\sum_h(n-a_h)a_h}\prod_h g(1/x;n-a_{h-1},n-a_h)\\
&=x^{\sum_h(n-a_h)a_h}\prod_h g(1/x;a_h,a_{h-1})\\
&=x^{\sum_h(n-a_h)a_{h-1}}\prod_h g(x;n-a_{h-1},n-a_h)\\
&=x^{\sum_h(n-a_h)a_{h-1}}\prod_h g(x;a_h,a_{h-1}).
\end{align*}

Ist von den Classen
\[
 \alpha=\Cl(a'_1\mid\cdots\mid a'_n),\qquad
 \beta=\Cl(b'_1\mid\cdots\mid b'_n)=\beta_0
\]
die erste durch die zweite theilbar, so ist (§ 8, III)
\begin{equation}
 t(\alpha\mid\beta)=
 {\psi(\beta)\psi(\beta_1)\cdots\psi(\beta_{n-2})
  \over
  \psi(\bar\beta)\psi(\bar\beta_1)\cdots\psi(\bar\beta_{n-2})}
 = {\psi(\beta)\psi(\beta_1)\cdots\psi(\beta_{n-1})
  \over
  \psi(\bar\beta)\psi(\bar\beta_1)\cdots\psi(\bar\beta_{n-1})}.
\tag{2}
\end{equation}
Dabei ergiebt sich die Bedeutung der Classen $\bar\beta=\bar\beta_0,\beta_i,\bar\beta_i$ $(i=1,\ldots,n-1)$ aus § 8, (5). Wir bezeichnen mit
\[
 a_h,
 \quad b_h=b_h^{(0)},\quad \bar b_h=\bar b_h^{(0)},
 \quad b_h^{(i)},\quad \bar b_h^{(i)}\qquad(i=1,\ldots,n-1)
\]
die Indices der Classen $\alpha,\beta=\beta_0,\bar\beta=\bar\beta_0,\beta_i,\bar\beta_i$ und erhalten, weil $\beta_i$ und $\bar\beta_i$ $(i=0,\ldots,n-1)$ vom Grade $n-i$ sind,
\begin{align}
 t(\alpha\mid\beta)
&=x^{\sum_h\sum_{i=0}^{n-1}\{(n-i-b_h^{(i)})b_h^{(i)}-(n-i-\bar b_h^{(i)})\bar b_h^{(i)}\}}
 \prod_h\prod_{i=0}^{n-1}
 {g(1/x;b_h^{(i)},b_{h-1}^{(i)})\over
 g(1/x;\bar b_h^{(i)},\bar b_{h-1}^{(i)})}.
\tag{3}
\end{align}
Wir wollen den eben gefundenen Ausdruck so umformen, dass er nur die Indices von $\alpha$ und $\beta$ enthält.

Für $a_h<i$ ist der Exponent von $\alpha$ kleiner als $h$, daher werden, wie aus § 8, (5) zu ersehen,
\[
 b_h^{(i)}=0,
 \qquad \bar b_h^{(i)}=0.
\]
Für $a_h=i$ ist $a'_i<h\le a'_{i+1}$, also
\[
 \bar b_h^{(i)}=0,
 \qquad b_h^{(i)}=b_h-i=b_h-a_h.
\]
Für $a_h>i$ ist $a'_{i+1}<h$, daher wird
\[
 b_h^{(i)}=\bar b_h^{(i)}=b_h-i.
\]
Hieraus folgt
\begin{equation}
 \sum_h\sum_{i=0}^{n-1}\{(n-i-b_h^{(i)})b_h^{(i)}-(n-i-\bar b_h^{(i)})\bar b_h^{(i)}\}
 =\sum_h(n-b_h)(b_h-a_h),
\tag{4}
\end{equation}
ferner
\[
\prod_{i=0}^{n-1}
 {g(1/x;b_h^{(i)},b_{h-1}^{(i)})\over
 g(1/x;\bar b_h^{(i)},\bar b_{h-1}^{(i)})}
 =g(1/x;b_h-a_{h-1},b_{h-1}-a_{h-1}).
\]
Man erhält also
\begin{align*}
\text{II.}\quad
 t(\alpha\mid\beta)
 &=x^{\sum_h(n-b_h)(b_h-a_h)}
 \prod_h g(1/x;b_h-a_{h-1},b_{h-1}-a_{h-1})\\
 &=x^{\sum_h(n-b_h)(b_{h-1}-a_{h-1})}
 \prod_h g(x;b_h-a_{h-1},b_{h-1}-a_{h-1})\\
 &=x^{\sum_h(n-b_h)(b_h-a_h)}
 \prod_h g(1/x;b_h-a_{h-1},b_h-b_{h-1})\\
 &=x^{\sum_h(n-b_h)(b_{h-1}-a_{h-1})}
 \prod_h g(x;b_h-a_{h-1},b_h-b_{h-1}).
\end{align*}
Nach § 7, I hat man jetzt
\[
 v(\alpha\mid\beta)={\psi(\alpha)\over\psi(\beta)}t(\alpha\mid\beta).
\]
Es ist aber
\[
 {\psi(\alpha)\over\psi(\beta)}=
 x^{\sum_h(n-a_h)a_h-\sum_h(n-b_h)b_h}
 \prod_h {g(1/x;a_h,a_{h-1})\over g(1/x;b_h,b_{h-1})},
\]
und mit den Formeln § 10, 3, 8), 11) geht dies in
\begin{align*}
\text{III.}\quad
 v(\alpha\mid\beta)&=x^{\sum_h(b_h-a_h)a_h}
 \prod_h g(1/x;b_h-a_{h-1},b_h-a_h)\\
 &=x^{\sum_h(b_h-a_h)a_{h-1}}
 \prod_h g(x;b_h-a_{h-1},b_h-a_h)\\
 &=x^{\sum_h(b_h-a_h)a_h}
 \prod_h g(1/x;b_h-a_{h-1},a_h-a_{h-1})\\
 &=x^{\sum_h(b_h-a_h)a_{h-1}}
 \prod_h g(x;b_h-a_{h-1},a_h-a_{h-1}).
\end{align*}
Alle diese Ausdrücke sind unter der Voraussetzung hergeleitet worden, dass $\alpha\ge\beta$ ist. Ist aber diese Voraussetzung nicht erfüllt, so kann man $h$ so wählen, dass $b_{h-1}<a_{h-1}\le b_h$ ist; dann verschwinden die vier unter II angegebenen Ausdrücke identisch, und dasselbe gilt von den vier Ausdrücken III. Da nun in dem hier besprochenen Fall thatsächlich $t(\alpha\mid\beta)=0$, $v(\alpha\mid\beta)=0$ ist, so bestehen die angegebenen Formeln ohne jede Einschränkung.

\section*{§ 12. Die Grade der Functionen $\psi$, $t$, $v$, $\bar\gamma$, $\bar\delta$.}

Da es im Folgenden nirgends mehr nöthig sein wird, Classen verschiedenen Grades neben einander zu betrachten, so sei jetzt festgesetzt, dass der Grad aller vorkommenden Classen mit $n$ bezeichnet werde.

Die Formeln § 11, I zeigen, dass $\psi(\alpha)$ eine ganze Function von $x$ mit ganzzahligen Coefficienten ist, dass der Coefficient der höchsten Potenz von $x$ gleich 1 und der Grad von $\psi(\alpha)$ gleich
\begin{equation}
 \sum_h(n-a_h)a_h
\tag{1}
\end{equation}
ist, wenn mit $a_h$ die Indices von $\alpha$ bezeichnet werden. Ist $\beta$ ein Divisor von $\alpha$, so sind (§ 11, II, III) $t(\alpha\mid\beta)$, $v(\alpha\mid\beta)$ ebenfalls ganze Functionen von $x$ mit ganzzahligen Coefficienten, und es ist auch bei ihnen der Coefficient der höchsten Potenz von $x$ gleich 1. Als die Grade der Functionen $t(\alpha\mid\beta)$, $v(\alpha\mid\beta)$ ergeben sich, wenn mit $b_h$ die Indices von $\beta$ bezeichnet werden, die Ausdrücke
\begin{equation}
 \sum_h(n-b_h)(b_h-a_h),
\tag{2}
\end{equation}
\begin{equation}
 \sum_h(b_h-a_h)a_h.
\tag{3}
\end{equation}
Für die Functionen $\bar\gamma(\gamma;\alpha_1,\ldots,\alpha_s)$, $\bar\delta(\delta;\alpha_1,\ldots,\alpha_s)$ haben wir nach § 9, IIa, IIIa die Relationen
\begin{equation}
\left\{
\begin{aligned}
 \bar\gamma(\gamma;\alpha_1,\ldots,\alpha_s)
 &=t(\gamma\mid\alpha_1)\cdots t(\gamma\mid\alpha_s)
 -\sum_{\gamma'\le\mu<\gamma}t(\gamma\mid\mu)\bar\gamma(\mu;\alpha_1,\ldots,\alpha_s),\\
 \bar\delta(\delta;\alpha_1,\ldots,\alpha_s)
 &=v(\alpha_1\mid\delta)\cdots v(\alpha_s\mid\delta)
 -\sum_{\delta'\ge\mu>\delta}v(\mu\mid\delta)\bar\delta(\mu;\alpha_1,\ldots,\alpha_s),
\end{aligned}
\right.
\tag{4}
\end{equation}
worin $\gamma'=[\alpha_1,\ldots,\alpha_s]$, $\delta'=(\alpha_1,\ldots,\alpha_s)$ ist und die Summation über alle diejenigen Classen $\mu$ zu erstrecken ist, welche den unter den Summenzeichen angegebenen Bedingungen genügen. Aus (4) und den Eigenschaften der Functionen $t$ und $v$ ersieht man, dass $\bar\gamma$ und $\bar\delta$ ganze Functionen von $x$ mit ganzzahligen Coefficienten sind, ferner, dass $\bar\gamma$ verschwindet, wenn die Bedingung $\gamma\ge\gamma'$ nicht erfüllt ist, dass $\bar\delta$ verschwindet, wenn die Bedingung $\delta\le\delta'$ nicht erfüllt ist.

Wir wollen zunächst den Grad von $\bar\delta$ näher untersuchen und setzen $\delta\le\delta'$ voraus. Der Coefficient der höchsten Potenz von $x$ kann in keiner $\bar\delta$-Function negativ sein, weil die Function für keinen Primzahlwerth von $x$ negativ sein kann, indem sie alsdann eine gewisse Anzahl darstellt. Aus (4) folgt daher, dass der Grad von $\bar\delta(\delta;\alpha_1,\ldots,
\alpha_s)$ höchstens gleich dem Grade von $v(\alpha_1\mid\delta)\cdots v(\alpha_s\mid\delta)$, d. i., wenn mit $a_h^{(i)}$, $d_h$ die Indices von $\alpha_i$ $(i=1,
\ldots,s)$, $\delta$ bezeichnet werden,
\begin{equation}
 \sum_{i=1}^s\sum_h(d_h-a_h^{(i)})a_h^{(i)}
\tag{5}
\end{equation}
sein kann und dass nur die beiden folgenden Fälle möglich sind: Entweder sind die Grade der Functionen
\begin{equation}
 v(\mu\mid\delta)\bar\delta(\mu;\alpha_1,\ldots,\alpha_s)
\tag{6}
\end{equation}
$(\delta'\ge\mu>\delta)$ sämmtlich kleiner als der Ausdruck (5); dann erreicht $\bar\delta$ den angegebenen Maximalgrad und der Coefficient der höchsten Potenz ist 1. Oder es steigt von den Functionen (6) eine bis zum Grade (5) an, dann bleibt der Grad von $\bar\delta$ hinter diesem Ausdruck zurück.

Wir wollen beweisen, dass die Function $\bar\delta(\delta;\alpha_1,
\ldots,\alpha_s)$ den angegebenen Maximalgrad (5) dann und nur dann erreicht, wenn die Classe $\delta$, welche bereits als Divisor von $\delta'$ vorausgesetzt wurde, ein Vielfaches des Grenztheilers $\delta''$ der Classen $\alpha_1,
\ldots,\alpha_s$ ist. Bezeichnen wir zu diesem Zweck mit $m_h$ die Indices der variablen Classe $\mu$, so ergiebt sich als Maximum für den Grad von
\[
 v(\mu\mid\delta)\bar\delta(\mu;\alpha_1,\ldots,\alpha_s)
\]
der Ausdruck
\begin{align}
&\sum_h(d_h-m_h)m_h+\sum_{i=1}^s\sum_h(m_h-a_h^{(i)})a_h^{(i)} \notag\\
&\qquad =\sum_{i=1}^s\sum_h(d_h-a_h^{(i)})a_h^{(i)}
 -\sum_h\biggl\{(d_h-m_h)\biggl(\sum_{i=1}^s a_h^{(i)}-m_h\biggr)\biggr\}.
\tag{7}
\end{align}
Es wird ferner (§ 10, 2, 7) u. 5)) $\delta$ dann und nur dann ein Vielfaches von $\delta''$ sein, wenn für jedes $h$
\begin{equation}
 d_h\le\sum_{i=1}^s a_h^{(i)}
\tag{8}
\end{equation}
ist. Ist nun diese Ungleichung beständig erfüllt, so folgt aus $\mu>\delta$, dass kein Glied der Summe
\begin{equation}
 \sum_h\biggl\{(d_h-m_h)\biggl(\sum_{i=1}^s a_h^{(i)}-m_h\biggr)\biggr\}
\tag{9}
\end{equation}
negativ wird. Dagegen muss wenigstens ein positives Glied in der Summe vorkommen, da die Gleichheit von $\mu$ und $\delta$ ausgeschlossen ist, also wenigstens für ein $h$ die Differenz $d_h-m_h$ und somit auch
\[
 \biggl(\sum_{i=1}^s a_h^{(i)}\biggr)-m_h>0
\]
ausfällt. In dem betrachteten Falle ist daher der Ausdruck (7) für jedes $\mu$ kleiner als der Ausdruck (5).

Wenn dagegen die Bedingungen (8) nicht durchweg erfüllt sind, so setzen wir
\[
 [\delta,\delta'']=\mu',
\]
dann ist der Index $m'_h$ von $\mu'$ gleich der kleineren der beiden Zahlen $\sum_{i=1}^s a_h^{(i)}$ und $d_h$, und die Classe $\mu'$ kommt, weil $\delta'\ge\mu'>\delta''$, unter den Classen $\mu$ vor. Aus $m'_h\le\sum_{i=1}^s a_h^{(i)}$ folgt nach dem soeben Bewiesenen, dass die Function $\bar\delta(\mu';\alpha_1,
\ldots,\alpha_s)$ wirklich bis zum Grade
\[
 \sum_{i=1}^s\sum_h(m'_h-a_h^{(i)})a_h^{(i)}
\]
ansteigt, die Function
\[
 v(\mu'\mid\delta)\bar\delta(\mu';\alpha_1,\ldots,\alpha_s)
\]
bis zum Grade
\[
 \sum_{i=1}^s\sum_h(d_h-a_h^{(i)})a_h^{(i)}
 -\sum_h\biggl\{(d_h-m'_h)\biggl(\sum_{i=1}^s a_h^{(i)}-m'_h\biggr)\biggr\}
\]
ansteigt, und da überdies das Product
\[
 (d_h-m'_h)\biggl(\sum_{i=1}^s a_h^{(i)}-m'_h\biggr)
\]
für jedes $h$ verschwindet, so wird der Grad von
\[
 v(\mu'\mid\delta)\bar\delta(\mu';\alpha_1,
\ldots,\alpha_s)
\]
gleich
\[
 \sum_{i=1}^s\sum_h(d_h-a_h^{(i)})a_h^{(i)}.
\]

In ganz ähnlicher Weise ergeben sich für die Function $\bar\gamma(\gamma;\alpha_1,
\ldots,\alpha_s)$ die folgenden Resultate: Ist $\gamma$ ein Vielfaches von $\gamma'=[\alpha_1,
\ldots,\alpha_s]$, so ist $\bar\gamma(\gamma;\alpha_1,
\ldots,\alpha_s)$ höchstens vom Grade
\begin{equation}
 \sum_{i=1}^s\sum_h(n-a_h^{(i)})(a_h^{(i)}-c_h),
\tag{10}
\end{equation}
wo mit $c_h$ die Indices von $\gamma$ bezeichnet worden sind. Die Function erreicht diesen Grad, wenn $\gamma$ Divisor von $\gamma''=[|\alpha_1,
\ldots,\alpha_s|]$ ist, und der Coefficient der höchsten Potenz von $x$ ist alsdann gleich 1; in den anderen Fällen wird der Maximalgrad nicht erreicht.

Da die Functionen $\bar\gamma$ und $\bar\delta$ in den Fällen, in welchen der angegebene Maximalgrad erreicht wird, nicht identisch verschwinden, so müssen sie, sobald $x$ eine gewisse positive Grösse überschreitet, stets positiv bleiben. Es wird sich später zeigen, dass dies schon für $x>1$ eintritt, dass daher die Functionen für keinen Primzahlwerth von $x$ verschwinden. Ferner wird es sich herausstellen, dass in den übrigen Fällen, für welche bisher nur gezeigt wurde, dass jener Maximalgrad nicht erreicht wird, die Functionen $\bar\gamma$ und $\bar\delta$ identisch verschwinden.

Damit ist alsdann der Beweis des Satzes § 7, V vollständig geführt. Um aber dahin zu gelangen, genügt es nicht, die Grade der Functionen $\bar\gamma$ und $\bar\delta$ zu betrachten. Wir müssen ihre Natur näher untersuchen und wollen zunächst diejenigen Functionen ins Auge fassen, aus denen sich, wie später nachgewiesen wird, die Functionen $\bar\gamma$ und $\bar\delta$ durch Multiplication unter Hinzufügung eines Factors, der nur eine Potenz von $x$ ist, herstellen lassen, in derselben Weise, wie die Functionen $\psi,t,v$ mit Hülfe der $g$-Functionen ausgedrückt wurden.

\section*{§ 13. Die $\omega$-Functionen.}

1. Wir bezeichnen mit
\[
 \omega(x\mid k\mid e\mid q_1,r_1;q_2,r_2;\ldots;q_s,r_s)
\]
die Summe
\begin{equation}
 \sum_m(-1)^m x^{{m(m-1)\over2}-km}
 g(x;e,m)g(x;q_1-m,r_1)g(x;q_2-m,r_2)\cdots g(x;q_s-m,r_s).
\tag{1}
\end{equation}
$x$ heisst das Argument, $k$ der erste, $e$ der zweite Parameter der $\omega$-Functionen, $q_1$ und $r_1$ bilden das erste, $q_2$ und $r_2$ das zweite, $\ldots$, $q_s$ und $r_s$ das $s$te Parameterpaar. Alle Parameter sollen ganze Zahlen sein, und die Summe erstreckt sich über alle ganzen Zahlen $m$. Ist $e\ge0$, und auf diesen Fall können wir uns hier beschränken, so ist nur eine endliche Anzahl von Gliedern der Summe von 0 verschieden. Die $\omega$-Function ist dann eine rationale Function von $x$, welche nur für $x=0$ und $x=\infty$ möglicherweise unendlich werden kann.

Die Function (1) soll eine ,,eigentliche`` $\omega$-Function heissen, wenn kein Parameter negativ und ferner
\[
 q_i\ge \begin{cases}e,\\ r_i,
 \end{cases}\qquad(i=1,\ldots,s)
\]
ist. In diesem Falle heisst der Ausdruck
\[
 k+r_1+\cdots+r_s-e
\]
die kritische Zahl der $\omega$-Function; ihre Bedeutung werden wir bald kennen lernen.

2. Wir behandeln zunächst die $\omega$-Function ohne Parameterpaare
\[
 \omega(x\mid k\mid e)=\sum_m(-1)^m x^{{m(m-1)\over2}-km}g(x;e,m).
\]
Es ist
\begin{equation}
 \omega(x\mid k\mid0)=1
\tag{2}
\end{equation}
und für $e>0$ folgt aus
\[
 g(x;e,m)=g(x;e-1,m-1)+x^m g(x;e-1,m)
\]
die Formel
\begin{align}
\omega(x\mid k\mid e)
&=\sum_m(-1)^m x^{{m(m-1)\over2}-km}g(x;e-1,m-1)\notag\\
&\quad+\sum_m(-1)^m x^{{m(m-1)\over2}-(k-1)m}g(x;e-1,m).
\tag{3}
\end{align}
Da nun
\[
 {m(m-1)\over2}-km={(m-1)(m-2)\over2}-(k-1)(m-1)-k,
\]
also
\begin{align*}
&\sum_m(-1)^m x^{{m(m-1)\over2}-km}g(x;e-1,m-1)\\
&\quad=-x^{-k}\sum_m(-1)^{m-1}x^{{(m-1)(m-2)\over2}-(k-1)(m-1)}g(x;e-1,m-1)\\
&\quad=-x^{-k}\sum_m(-1)^m x^{{m(m-1)\over2}-(k-1)m}g(x;e-1,m),
\end{align*}
ist, so erhält man aus (3)
\begin{equation}
 \omega(x\mid k\mid e)=(1-x^{-k})\omega(x\mid k-1\mid e-1)
\tag{4}
\end{equation}
und weiter, unter Berücksichtigung von (2),
\begin{equation}
 \omega(x\mid k\mid e)=(1-x^{-k})(1-x^{-k+1})\cdots(1-x^{-k+e-1}).
\tag{5}
\end{equation}
Die betrachtete $\omega$-Function sei nun eine eigentliche, also
\[
 k\ge0,
 \qquad e\ge0.
\]
Dann ist $k-e$ die kritische Zahl und man sieht aus (5), dass die Function identisch verschwindet, wenn die kritische Zahl negativ ist, dass sie dagegen anderenfalls für jedes $x>1$ einen positiven von 0 verschiedenen Werth hat.

3. Für die $\omega$-Function mit $s$ Parameterpaaren, von welcher wir zunächst nur voraussetzen wollen, dass ihr zweiter Parameter $e$ nicht negativ ist, erhält man durch Anwendung der Relation
\[
 g(x;q_s-m,r_s)=g(x;q_s-1-m,r_s-1)+x^{r_s}g(x;q_s-1-m,r_s)
\]
die Formel
\begin{align}
&\omega(x\mid k\mid e\mid q_1,r_1;\ldots;q_s,r_s)\notag\\
&\quad=x^{r_s}\omega(x\mid k\mid e\mid q_1,r_1;\ldots;q_{s-1},r_{s-1};q_s-1,r_s)\notag\\
&\qquad+\omega(x\mid k\mid e\mid q_1,r_1;\ldots;q_{s-1},r_{s-1};q_s-1,r_s-1),
\tag{6}
\end{align}
und hieraus allgemeiner
\begin{align}
&\omega(x\mid k\mid e\mid q_1,r_1;\ldots;q_s,r_s)\notag\\
&\quad=\sum_{l=0}^t C'_l\,
 \omega(x\mid k\mid e\mid q_1,r_1;\ldots;q_{s-1},r_{s-1};q_s-t,r_s-l),
\tag{7}
\end{align}
wo $t$ irgend eine ganze nicht negative Zahl bezeichnet. Der Coefficient $C'_l$ ist, wie man durch Induction leicht findet, gleich
\[
 x^{l(r_s-t)}g(x;t,l),
\]
das, worauf es uns ankommt, ist, dass er für jeden positiven Werth von $x$ selbst positiv ist.

4. Wir betrachten den besonderen Fall, dass $q_s=e$ ist, und setzen $e\ge0$ voraus. Aus
\[
 g(x;e,m)g(x;e-m,r_s)=g(x;e,r_s)g(x;e-r_s,m)
\]
ergiebt sich:
\begin{align}
&\omega(x\mid k\mid e\mid q_1,r_1;\ldots;q_{s-1},r_{s-1};e,r_s)\notag\\
&\quad=g(x;e,r_s)\omega(x\mid k\mid e-r_s\mid q_1,r_1;\ldots;q_{s-1},r_{s-1}).
\tag{8}
\end{align}
Die $\omega$-Function auf der rechten Seite enthält ein Parameterpaar weniger. Der Factor $g(x;e,r_s)$ verschwindet, wenn $r_s<0$ oder $r_s>e$ ist; sonst hat er für positives $x$ einen positiven Werth.

5. Es sei jetzt
\[
 \omega(x\mid k\mid e\mid q_1,r_1;\ldots;q_s,r_s)
\]
eine eigentliche $\omega$-Function. Dann ist $q_s-e\ge0$, und wenn wir Formel (7) für $t=q_s-e$ und darauf Formel (8) zur Anwendung bringen, so erhalten wir
\begin{align}
&\omega(x\mid k\mid e\mid q_1,r_1;\ldots;q_s,r_s)\notag\\
&\quad=\sum_{l=0}^{q_s-e}C'_l g(x;e,r_s-l)\,
 \omega(x\mid k\mid e-r_s+l\mid q_1,r_1;\ldots;q_{s-1},r_{s-1}).
\tag{9}
\end{align}
Setzt man $r_s-l=h$, so wird demnach
\begin{align}
&\omega(x\mid k\mid e\mid q_1,r_1;\ldots;q_s,r_s)\notag\\
&=\sum_{\substack{r_s+e-q_s\le h\le r_s\\0\le h\le e}}
 C'_{r_s-h}g(x;e,h)\omega(x\mid k\mid e-h\mid q_1,r_1;\ldots;q_{s-1},r_{s-1})\notag\\
&=\sum_{\substack{r_s+e-q_s\le h\le r_s\\0\le h\le e}}
 C_h\,\omega(x\mid k\mid e-h\mid q_1,r_1;\ldots;q_{s-1},r_{s-1}).
\tag{10}
\end{align}
Die Summe erstreckt sich über alle $h$, welche den Bedingungen
\[
 \max(0,r_s+e-q_s)\le h\le \min(e,r_s)
\]
genügen, und die Coefficienten $C$ haben für jedes positive $x$ einen positiven Werth. Die $\omega$-Functionen auf der rechten Seite von (10) sind eigentliche und enthalten nur $s-1$ Parameterpaare. Es ist daraus zu ersehen, dass man jede eigentliche $\omega$-Function linear und homogen durch eigentliche $\omega$-Functionen ohne Parameterpaare so ausdrücken kann, dass alle Coefficienten für positives $x$ positiv werden.

Die kritische Zahl der betrachteten $\omega$-Function ist
\[
 k+r_1+\cdots+r_s-e,
\]
die Function
\[
 \omega(x\mid k\mid e-h\mid q_1,r_1;\ldots;q_{s-1},r_{s-1})
\]
hat die kritische Zahl
\[
 k+r_1+\cdots+r_{s-1}+h-e,
\]
welche für alle in Betracht kommenden Werthe von $h$
\[
 \le k+r_1+\cdots+r_s-e
\]
ist. Daraus folgt, dass eine eigentliche $\omega$-Function mit negativer kritischer Zahl sich linear und homogen durch ebensolche $\omega$-Functionen ohne Parameterpaare ausdrücken lässt und mithin verschwindet.

Wir wollen nun annehmen, dass die kritische Zahl
\[
 k+r_1+\cdots+r_s-e\ge0
\]
sei. Auf der rechten Seite von (10) verschwinden die $\omega$-Functionen mit negativer kritischer Zahl, es bleibt aber wenigstens noch eine $\omega$-Function zurück, deren kritische Zahl nicht negativ ist, diejenige nämlich, welche man erhält, wenn man $h$ gleich der kleineren der beiden Zahlen $e$ und $r_s$ setzt. Man erkennt hieraus, dass man in diesem Falle die $\omega$-Function durch eigentliche $\omega$-Functionen ohne Parameterpaare mit nicht negativer kritischer Zahl linear und homogen ausdrücken kann, und dass in diesem Ausdruck wenigstens eine derartige Function wirklich vorkommt. Da ferner die auftretenden Coefficienten für positives $x$ selbst positiv sind, die $\omega$-Functionen aber, sobald $x>1$ ist, positiv sind, so hat auch $\omega(x\mid k\mid e\mid q_1,r_1;\ldots;q_s,r_s)$ für $x>1$ einen positiven Werth.

Damit sind wir zu dem wichtigen Resultat gelangt:

I. Eine eigentliche $\omega$-Function, deren kritische Zahl $<0$ ist, verschwindet identisch.

II. Eine eigentliche $\omega$-Function, deren kritische Zahl $\ge0$ ist, hat für $x>1$ einen positiven Werth.

\section*{§ 14. Die Functionen $\bar t$ und $\bar v$.}

Für $e>0$ ist
\begin{equation}
 0=\omega(x\mid0\mid e)=\sum_m(-1)^m x^{m(m-1)/2}g(x;e,m).
\tag{1}
\end{equation}
Sind daher $a,b,c$ ganze Zahlen, welche den Bedingungen
\begin{equation}
 0\le c\le a<b
\tag{2}
\end{equation}
genügen, so geht, wenn man $e$ durch $b-a$ und $m$ durch $b-m$ ersetzt, Gleichung (1) über in
\[
 \sum_m(-1)^{b-m}x^{(b-m)(b-m-1)/2}g(x;b-a,b-m)=0,
\]
und man erhält durch Multiplication mit
\[
 x^{c(b-a)}g(x;b-c,a-c)
\]
und unter Berücksichtigung der Identitäten
\begin{align*}
 g(x;b-a,b-m)g(x;b-c,a-c)
 &=g(x;m-c,a-c)g(x;b-c,b-m)\\
 &=g(x;m-c,m-a)g(x;b-c,b-m)
\end{align*}
die Relation
\begin{equation}
 \sum_m(-1)^{b-m}x^{(b-m)(b-m-1)/2+c(b-a)}
 g(x;m-c,m-a)g(x;b-c,b-m)=0.
\tag{3}
\end{equation}
Es seien nun $\alpha$ und $\beta$ zwei Classen, deren Indices mit $a_h,b_h$ bezeichnet werden mögen. Die Summe
\begin{align}
F_h&=\sum_{m_h}(-1)^{b_h-m_h}
 x^{(b_h-m_h)(b_h-m_h-1)/2+(m_h-a_h)a_{h-1}+(b_h-m_h)b_{h-1}}\notag\\
&\quad\cdot g(x;m_h-a_{h-1},m_h-a_h)g(x;b_h-b_{h-1},b_h-m_h)
\tag{4}
\end{align}
hat dann für jedes $h$, wenn $m_h$ alle ganzen Zahlen durchläuft, nur eine endliche Anzahl nicht verschwindender Glieder. Man sieht ferner, dass $F_h=1$ wird, wenn $h\le0$ oder grösser ist als alle Exponenten von $\alpha$ und $\beta$.

Da man sich bei der Summation auf diejenigen Werthe von $m_h$ beschränken kann, welche den Bedingungen
\begin{equation}
 b_h\ge m_h\ge\begin{cases}b_{h-1},\\ a_h
 \end{cases}
\tag{5}
\end{equation}
genügen, so verschwindet $F_h$ für $a_h>b_h$. $F_h$ verschwindet aber auch, wenn $a_{h-1}=b_{h-1}$ und zugleich $a_h<b_h$ ist. Denn alsdann geht, wenn man $a_{h-1}=b_{h-1}=c$, $a_h=a$, $b_h=b$ setzt, der Ausdruck (4) in (3) über und die Bedingung (2) ist erfüllt. Ist dagegen $a_{h-1}=b_{h-1}$, $a_h=b_h$, so ergiebt sich $F_h=1$. Das Product
\[
 \prod_hF_h
\]
hat daher den Werth 1 oder 0, je nachdem die Classen $\alpha$ und $\beta$ gleich oder verschieden sind. Aus der in § 7 für die Functionen $\bar\gamma$ und $\bar\delta$ gegebenen Definition folgt für den Fall, dass die Anzahl $s$ der in $\bar\gamma(\gamma;\alpha_1,\ldots,\alpha_s)$, $\bar\delta(\delta;\alpha_1,
\ldots,\alpha_s)$ vorkommenden Classen $\alpha$ sich auf 1 reducirt, $\bar\gamma(\gamma;\alpha)=1$ oder $0$, $\bar\delta(\delta;\alpha)=1$ oder $0$, je nachdem $\gamma$ und $\alpha$, bez. $\delta$ und $\alpha$ gleich oder verschieden sind. Man kann daher das für $\prod_hF_h$ erhaltene Resultat ausdrücken durch die Gleichung
\begin{equation}
 \prod_hF_h=\bar\delta(\alpha;\beta).
\tag{6}
\end{equation}
Aus (4) folgt
\begin{align}
\prod_hF_h
&=\sum_{\ldots,m_0,m_1,\ldots}
\biggl\{(-1)^{\sum_h(b_h-m_h)}
 x^{\sum_h{(b_h-m_h)(b_h-m_h-1)\over2}+\sum_h(b_h-m_h)b_{h-1}}
 \prod_h g(x;b_h-b_{h-1},b_h-m_h)\notag\\
&\qquad\qquad\cdot
 x^{\sum_h(m_h-a_h)a_{h-1}}
 \prod_h g(x;m_h-a_{h-1},m_h-a_h)\biggr\}.
\tag{7}
\end{align}
Man erhält ein Glied dieser Summe, indem man dem $m_h$ für jedes $h$ einen bestimmten Werth nach Belieben beilegt. Man darf sich aber auf solche Werthenreihen für $m_h$ beschränken, welche den Bedingungen (5) genügen; andere Werthenreihen liefern verschwindende Glieder und dürfen also nach Willkür hinzugenommen oder fortgelassen werden. Man erhält daher die Summe (7) vollständig, wenn man alle diejenigen Werthenreihen $m_h$ berücksichtigt, welche den Bedingungen
\begin{equation}
 b_h\ge m_h\ge a_h,
 \qquad m_h\ge m_{h-1},
\tag{8}
\end{equation}
oder gar alle diejenigen, welche den Bedingungen
\begin{equation}
 m_h\ge m_{h-1}\ge0,
 \qquad \lim_{h=\infty}m_h=n,
\tag{9}
\end{equation}
genügen. Die letzteren besagen genau, dass die $m_h$ die Indices einer Classe $\mu$ vom Grade $n$ bilden, während durch (8) noch ausgedrückt wird, dass diese Classe die Ungleichung
\[
 \alpha\ge\mu\ge\beta
\]
erfüllt. Unter der Voraussetzung, dass die $m_h$ die Indices einer Classe $\mu$ bilden, ist der zweite Factor unter dem Summenzeichen in (7) nichts anderes als unsere Function $v(\alpha\mid\mu)$, während der erste eine wohldefinirte Function von $\beta$ und $\mu$ darstellt, die mit $\bar v(\mu\mid\beta)$ bezeichnet werden mag:
\begin{equation}
 \bar v(\mu\mid\beta)=(-1)^{\sum_h(b_h-m_h)}
 \prod_h x^{(b_h-m_h)(b_h-m_h-1)/2+(b_h-m_h)b_{h-1}}
 g(x;b_h-b_{h-1},b_h-m_h).
\tag{10}
\end{equation}
Die Gleichung (6) geht hiernach über in
\begin{equation}
 \sum_{\mu}v(\alpha\mid\mu)\bar v(\mu\mid\beta)=\bar\delta(\alpha;\beta),
\tag{11}
\end{equation}
wo die Summe über alle Classen $\mu$ oder nur die zwischen $\alpha$ und $\beta$ gelegenen zu erstrecken ist.

Aus (11) folgt
\begin{align*}
&\sum_\nu\sum_\mu v(\alpha\mid\nu)\bar v(\nu\mid\mu)v(\mu\mid\beta)
 =\sum_\mu\bar\delta(\alpha;\mu)v(\mu\mid\beta),\\
&\sum_\nu\sum_\mu v(\alpha\mid\nu)\bar v(\nu\mid\mu)v(\mu\mid\beta)=v(\alpha\mid\beta),
\tag{12}
\end{align*}
ferner
\begin{equation}
 \bar v(\alpha\mid\beta)=\bar\delta(\alpha;\beta)
 -\sum_{\mu<\alpha}v(\alpha\mid\mu)\bar v(\mu\mid\beta).
\tag{13}
\end{equation}
Die Summe rechts erstreckt sich über die Classen $\mu$, welche $<\alpha$ sind; ist also für alle diese Classen $\bar v(\mu\mid\beta)$ bekannt, so lehrt (13), $\bar v(\alpha\mid\beta)$ zu berechnen. Demnach ist durch die Relation (11) die Function $\bar v(\alpha\mid\beta)$ vollkommen definirt.

Ist $\alpha$ kein Vielfaches von $\beta$, so ist $\bar v(\alpha\mid\beta)=0$; man ersieht dies aus (13), wenn man für $\mu<\alpha$ die Gleichung $\bar v(\mu\mid\beta)=0$ als bewiesen annimmt. Sodann zeigt (13), dass für $\alpha=\beta$ $\bar v(\alpha\mid\beta)=1$ wird. In beiden Fällen hat man:
\begin{equation}
 \sum_\mu\bar v(\alpha\mid\mu)v(\mu\mid\beta)=\bar\delta(\alpha;\beta).
\tag{14}
\end{equation}
Indem wir nun für den noch nicht betrachteten Fall $\alpha>\beta$ die Richtigkeit von (14) beweisen wollen, können wir wieder für jede Classe $\nu<\alpha$ die Gleichung
\[
 \sum_\mu\bar v(\nu\mid\mu)v(\mu\mid\beta)=\bar\delta(\nu;\beta)
\]
als bewiesen ansehen. Aus dieser folgt
\begin{align}
\sum_{\nu<\alpha}\sum_\mu v(\alpha\mid\nu)\bar v(\nu\mid\mu)v(\mu\mid\beta)
 &=\sum_{\nu<\alpha}v(\alpha\mid\nu)\bar\delta(\nu;\beta),\notag\\
\sum_{\nu<\alpha}\sum_\mu v(\alpha\mid\nu)\bar v(\nu\mid\mu)v(\mu\mid\beta)&=v(\alpha\mid\beta).
\tag{15}
\end{align}
Subtrahirt man (15) von (12), so erhält man rechts Null und links
\[
 \sum_\mu v(\alpha\mid\alpha)\bar v(\alpha\mid\mu)v(\mu\mid\beta)
 =\sum_\mu\bar v(\alpha\mid\mu)v(\mu\mid\beta).
\]
Damit ist (14) allgemein bewiesen.

Für die Function
\begin{equation}
 \bar t(\mu\mid\beta)={\psi(\beta)\over\psi(\mu)}\bar v(\mu\mid\beta)
\tag{16}
\end{equation}
folgen aus
\[
 t(\alpha\mid\mu)\bar t(\mu\mid\beta)
 ={\psi(\beta)\over\psi(\alpha)}v(\alpha\mid\mu)\bar v(\mu\mid\beta),
\]
\[
 \bar t(\alpha\mid\mu)t(\mu\mid\beta)
 ={\psi(\beta)\over\psi(\alpha)}\bar v(\alpha\mid\mu)v(\mu\mid\beta),
\]
\[
 {\psi(\beta)\over\psi(\alpha)}\bar\delta(\alpha;\beta)=\bar\delta(\alpha;\beta)=\bar\gamma(\alpha;\beta)
\]
und unter Berücksichtigung von (11), (14)
\begin{equation}
 \sum_\mu t(\alpha\mid\mu)\bar t(\mu\mid\beta)=\bar\gamma(\alpha;\beta),
\tag{17}
\end{equation}
\begin{equation}
 \sum_\mu\bar t(\alpha\mid\mu)t(\mu\mid\beta)=\bar\gamma(\alpha;\beta).
\tag{18}
\end{equation}
Die Ausrechnung von $\bar t$ mit Hülfe der für $v$ und $\psi$ geltenden Formeln ergiebt nach (16)\footnote{Vgl. die Berechnung von $v(\alpha\mid\beta)$, S. 41.}
\begin{equation}
 \bar t(\mu\mid\beta)=(-1)^{\sum_h(b_h-m_h)}
 \prod_h x^{(b_h-m_h)(b_h-m_h-1)/2+(n-m_h)(b_{h-1}-m_{h-1})}
 g(x;m_h-m_{h-1},b_{h-1}-m_{h-1}).
\tag{19}
\end{equation}
Nach § 9, I, II ist
\begin{equation}
 \sum_\mu t(\nu\mid\mu)\bar\gamma(\mu;\alpha_1,\ldots,\alpha_s)
 =t(\nu\mid\alpha_1)t(\nu\mid\alpha_2)\cdots t(\nu\mid\alpha_s),
\tag{20}
\end{equation}
\begin{equation}
 \sum_\mu v(\mu\mid\lambda)\bar\delta(\mu;\alpha_1,\ldots,\alpha_s)
 =v(\alpha_1\mid\lambda)v(\alpha_2\mid\lambda)\cdots v(\alpha_s\mid\lambda).
\tag{21}
\end{equation}
Multiplicirt man (20), (21) bez. mit $\bar t(\gamma\mid\nu)$, $\bar v(\lambda\mid\delta)$ und summirt man dann über $\nu$ bez. $\lambda$, so erhält man unter Anwendung von (18), (11) links
\[
 \sum_{\nu,\mu}\bar t(\gamma\mid\nu)t(\nu\mid\mu)\bar\gamma(\mu;\alpha_1,\ldots,\alpha_s)
 =\sum_\mu\bar\gamma(\gamma;\mu)\bar\gamma(\mu;\alpha_1,\ldots,\alpha_s)
 =\bar\gamma(\gamma;\alpha_1,\ldots,\alpha_s),
\]
bez.
\[
 \sum_{\lambda,\mu}v(\mu\mid\lambda)\bar v(\lambda\mid\delta)\bar\delta(\mu;\alpha_1,\ldots,\alpha_s)
 =\sum_\mu\bar\delta(\mu\mid\delta)\bar\delta(\mu;\alpha_1,\ldots,\alpha_s)
 =\bar\delta(\delta;\alpha_1,\ldots,\alpha_s),
\]
und somit die Gleichungen
\begin{equation}
 \bar\gamma(\gamma;\alpha_1,
\ldots,\alpha_s)=\sum_\nu \bar t(\gamma\mid\nu)t(\nu\mid\alpha_1)\cdots t(\nu\mid\alpha_s),
\tag{22}
\end{equation}
\begin{equation}
 \bar\delta(\delta;\alpha_1,
\ldots,\alpha_s)=\sum_\lambda \bar v(\lambda\mid\delta)v(\alpha_1\mid\lambda)\cdots v(\alpha_s\mid\lambda).
\tag{23}
\end{equation}

\section*{§ 15. Darstellung von $\bar\gamma$ und $\bar\delta$ durch $\omega$-Functionen.}

Die beiden letzten Gleichungen führen unmittelbar zur Darstellung von $\bar\gamma(\gamma;\alpha_1,
\ldots,\alpha_s)$, $\bar\delta(\delta;\alpha_1,
\ldots,\alpha_s)$ durch $\omega$-Functionen. Es seien mit $a_h^{(i)},c_h,d_h$ die Indices von $\alpha_i$ $(i=1,\ldots,s)$, $\gamma$, $\delta$ bezeichnet.

Wir wollen nur die Function $\bar\delta$ etwas eingehender behandeln. Man erhält, indem man die für die Functionen $v,\bar v$ gefundenen Ausdrücke substituirt, die Gleichung
\begin{align}
&\bar v(\mu\mid\delta)v(\alpha_1\mid\mu)\cdots v(\alpha_s\mid\mu)\notag\\
&=\prod_h\biggl\{(-1)^{d_h-m_h}
 x^{(d_h-m_h)(d_h-m_h-1)/2+(d_h-m_h)d_{h-1}+(m_h-a_h^{(1)})a_{h-1}^{(1)}+\cdots+(m_h-a_h^{(s)})a_{h-1}^{(s)}}
\notag\\
&\qquad\cdot g(x;d_h-d_{h-1},d_h-m_h)g(x;m_h-a_{h-1}^{(1)},a_h^{(1)}-a_{h-1}^{(1)})
\cdots g(x;m_h-a_{h-1}^{(s)},a_h^{(s)}-a_{h-1}^{(s)})\biggr\},
\tag{1}
\end{align}
wo die $m_h$ die Indices von $\mu$ bedeuten. Um $\bar\delta(\delta;\alpha_1,
\ldots,\alpha_s)$ zu erhalten, hat man die Summe aller solcher Producte zu bilden, welche man für die verschiedenen Werthenreihen $m_h$ erhält. Die $m_h$ unterliegen dabei den für die Indices einer Classe $n$ten Grades geltenden Beschränkungen. Da man aber sofort sieht, dass für jede Reihe von Zahlen $m_h$, welche nicht Indices einer Classe $n$ten Grades sind, das Product auf der rechten Seite von (1) verschwindet, so kann man die angegebene Beschränkung aufheben. Alsdann erhält man:
\begin{align}
&\bar\delta(\delta;\alpha_1,\ldots,\alpha_s)
 =\sum_\mu \bar v(\mu\mid\delta)v(\alpha_1\mid\mu)\cdots v(\alpha_s\mid\mu)\notag\\
&=\prod_h\biggl\{\sum_{m_h}\biggl[(-1)^{d_h-m_h}
 x^{(d_h-m_h)(d_h-m_h-1)/2+(d_h-m_h)d_{h-1}+(m_h-a_h^{(1)})a_{h-1}^{(1)}+\cdots+(m_h-a_h^{(s)})a_{h-1}^{(s)}}\notag\\
&\qquad\cdot g(x;d_h-d_{h-1},d_h-m_h)g(x;m_h-a_{h-1}^{(1)},a_h^{(1)}-a_{h-1}^{(1)})\cdots g(x;m_h-a_{h-1}^{(s)},a_h^{(s)}-a_{h-1}^{(s)})\biggr]\biggr\}.
\tag{2}
\end{align}
Setzt man zur Abkürzung
\begin{align}
(d_h-a_h^{(1)})a_{h-1}^{(1)}+\cdots+(d_h-a_h^{(s)})a_{h-1}^{(s)}&=f_h,
\tag{3}\\
a_{h-1}^{(1)}+\cdots+a_{h-1}^{(s)}-d_{h-1}&=k_{h-1},
\tag{4}\\
d_h-d_{h-1}&=e_h,
\tag{5}\\
d_h-a_{h-1}^{(i)}=q_h^{(i)},\qquad a_h^{(i)}-a_{h-1}^{(i)}&=r_h^{(i)}\quad(i=1,\ldots,s),
\tag{6}
\end{align}
so wird die in (2) vorkommende Summe gleich
\[
 x^{f_h}\omega(x\mid k_{h-1}\mid e_h\mid q_h^{(1)},r_h^{(1)};\ldots;q_h^{(s)},r_h^{(s)}),
\]
und man erhält
\[
\text{I.}\qquad
 \bar\delta(\delta;\alpha_1,
\ldots,\alpha_s)=\prod_h x^{f_h}\,
 \omega(x\mid k_{h-1}\mid e_h\mid q_h^{(1)},r_h^{(1)};\ldots;q_h^{(s)},r_h^{(s)}).
\]
Jede der hier auftretenden $\omega$-Functionen
\begin{equation}
 \omega_h=\omega(x\mid k_{h-1}\mid e_h\mid q_h^{(1)},r_h^{(1)};\ldots;q_h^{(s)},r_h^{(s)})
\tag{7}
\end{equation}
wird dargestellt durch eine Summe, welche nur eine endliche Anzahl nicht verschwindender Glieder enthält, weil nach (5) $e_h\ge0$ ist. Wenn $h\le0$ oder grösser ist als alle Exponenten der Classen $\delta,\alpha_1,\ldots,
\alpha_s$, so ist $f_h=0$ und $\omega_h=1$, weil alle Parameter von $\omega_h$ ausser dem ersten gleich 0 werden. In dem Product I sind daher nur eine endliche Anzahl von Factoren von 1 verschieden.

Nach (5) und (6) ist
\begin{equation}
 e_h\ge0,
 \qquad r_h^{(i)}\ge0.
\tag{8}
\end{equation}
Damit $\bar\delta(\delta;\alpha_1,
\ldots,\alpha_s)$ nicht identisch verschwinde, ist nothwendig und hinreichend, dass kein $\omega_h$ verschwindet. Wir wollen annehmen, es sei dies der Fall, und sehen, welche Schlüsse wir daraus ziehen können.

Für $h\le0$ ist $q_h^{(i)}=0=r_h^{(i)}$. Ist jetzt $h$ irgend eine ganze Zahl, für welche $q_{h-1}^{(i)}\ge r_{h-1}^{(i)}$ ist, so hat man
\[
 0\le q_{h-1}^{(i)}-r_{h-1}^{(i)}=d_{h-1}-a_{h-1}^{(i)}\le d_h-a_{h-1}^{(i)}=q_h^{(i)},
\]
also $q_h^{(i)}\ge0$. Wäre nun $q_h^{(i)}<r_h^{(i)}$, so würde aus $r_h^{(i)}>q_h^{(i)}\ge0$ sich $\omega_h=0$ ergeben. Dies widerspricht der Voraussetzung, und daher ist für jedes $h$
\begin{equation}
 q_h^{(i)}\ge r_h^{(i)}\qquad(i=1,
\ldots,s).
\tag{9}
\end{equation}
Aus $q_h^{(i)}-r_h^{(i)}=d_h-a_h^{(i)}=q_{h+1}^{(i)}-e_{h+1}$ folgt, dass die Bedingung (9) mit jeder der beiden folgenden äquivalent ist:
\begin{equation}
 d_h\ge a_h^{(i)}\qquad(i=1,
\ldots,s),
\tag{10}
\end{equation}
\begin{equation}
 q_h^{(i)}\ge e_h\qquad(i=1,
\ldots,s).
\tag{11}
\end{equation}
Für $h\le0$ ist $k_h=0$. Ist $h$ irgend eine ganze Zahl, für welche $k_{h-1}\ge0$ ist, so ist zufolge (8), (9), (10), (11) $\omega_h$ eine eigentliche $\omega$-Function. Ihre kritische Zahl ist
\[
 k_{h-1}+r_h^{(1)}+\cdots+r_h^{(s)}-e_h=k_h
\]
und muss, damit $\omega_h$ nicht verschwinde, $\ge0$ sein. Man hat also für jedes $h$
\begin{equation}
 k_h\ge0.
\tag{12}
\end{equation}
Umgekehrt ergiebt sich aus (8) bis (12), dass jedes $\omega_h$ eine eigentliche $\omega$-Function mit nicht negativer kritischer Zahl ist, und hieraus folgt (§ 13, II) nicht nur, dass $\bar\delta$ nicht identisch verschwindet, sondern auch, dass $\bar\delta$ für jedes $x>1$ einen positiven Werth hat. Die Relation (12) ist mit
\begin{equation}
 d_h\le\sum_{i=1}^s a_h^{(i)}
\tag{13}
\end{equation}
äquivalent, die Relation (10) mit (9) und (11), während (8) von selbst erfüllt ist. Die nothwendigen und hinreichenden Bedingungen für das Nichtverschwinden der Function
\[
 \bar\delta(\delta;\alpha_1,
\ldots,\alpha_s)
\]
werden also durch die Relationen (10) und (13) dargestellt; und von diesen besagt die erste, dass $\delta$ ein Divisor von $\delta'=(\alpha_1,
\ldots,\alpha_s)$, die zweite, dass $\delta$ ein Vielfaches von $\delta''=(|\alpha_1,
\ldots,\alpha_s|)$ sein muss.

Für die Function $\bar\gamma(\gamma;\alpha_1,
\ldots,\alpha_s)$ erhält man in derselben Weise, wenn man hier
\[
\begin{gathered}
 (n-a_{h+1}^{(1)})(a_h^{(1)}-c_h)+\cdots+(n-a_{h+1}^{(s)})(a_h^{(s)}-c_h)=f_h,\\
 (n-a_{h+1}^{(1)})+\cdots+(n-a_{h+1}^{(s)})-(n-c_{h+1})=k_{h+1},\\
 c_{h+1}-c_h=e_h,\\
 a_{h+1}^{(i)}-c_h=q_h^{(i)},
 \qquad a_{h+1}^{(i)}-a_h^{(i)}=r_h^{(i)}\quad(i=1,\ldots,s)
\end{gathered}
\]
setzt,
\[
\text{II.}\qquad
 \bar\gamma(\gamma;\alpha_1,
\ldots,\alpha_s)=\prod_h x^{f_h}\,
 \omega(x\mid k_{h+1}\mid e_h\mid q_h^{(1)},r_h^{(1)};\ldots;q_h^{(s)},r_h^{(s)}).
\]
Damit $\bar\gamma$ nicht verschwinde, müssen alle
\[
 \omega_h=\omega(x\mid k_{h+1}\mid e_h\mid q_h^{(1)},r_h^{(1)};\ldots;q_h^{(s)},r_h^{(s)})
\]
eigentliche $\omega$-Functionen mit nicht negativer kritischer Zahl werden, und es hat alsdann $\bar\gamma$ für jedes $x>1$ einen positiven Werth. Die nothwendigen und hinreichenden Bedingungen hierfür werden dargestellt durch die Relationen
\[
 c_h\le a_h^{(i)},
 \qquad
 c_h\ge n-\sum_{i=1}^s(n-a_h^{(i)}),
\]
von denen die erste besagt, dass $\gamma$ ein Vielfaches von $\gamma'=[\alpha_1,
\ldots,\alpha_s]$, die zweite, dass $\gamma$ ein Divisor von $\gamma''=[|\alpha_1,
\ldots,\alpha_s|]$ sein muss.

Hiermit ist Satz V, § 7 vollständig bewiesen.

Charlottenburg, den 9. April 1898.

\end{document}
